% ============================================================================
% SAID - CONTRATO DE COMUNICACION UART v2.0.0 (FINAL)
% ============================================================================
% Especificación técnica final y cerrada del protocolo de comunicación lógico
% entre el microcontrolador ESP32 y el microcontrolador Arduino Mega para la
% máquina expendedora SAID.
% ============================================================================

\documentclass[a4paper,11pt]{article}

% ============================================================================
% METADATA DEL DOCUMENTO
% ============================================================================
% Modificar estos comandos actualiza de forma global encabezados, pies de página,
% portada y referencias del documento sin repetir cadenas manualmente.

\newcommand{\DocumentTitle}{Contrato UART SAID}
\newcommand{\DocumentSubtitle}{Especificación lógica de comunicación ESP32--Arduino Mega}
\newcommand{\DocumentType}{Especificación de interfaz}
\newcommand{\DocumentId}{SAID-ARCH-UART}
\newcommand{\DocumentVersion}{2.0.0}
\newcommand{\DocumentStatus}{FINAL}
\newcommand{\DocumentDate}{12 de septiembre de 2026}
\newcommand{\DocumentOwner}{Subgrupo de Arquitectura}
\newcommand{\DocumentAuthors}{César Euresti, Diego Eduardo Zapata Aguilar, Dulce Yaretzi Martinez Perez, Myranda Belen Infante Castillo}
\newcommand{\DocumentConsumers}{Programación ESP32 / Programación Arduino Mega}

% ============================================================================
% PAQUETES Y CONFIGURACION DE IDIOMA
% ============================================================================
% Se cargan los paquetes estándar necesarios para tipografía técnica, tablas,
% bloques de código, gráficos TikZ y gestión de enlaces PDF.

\usepackage[utf8]{inputenc}
\usepackage[T1]{fontenc}
\usepackage[spanish,es-nodecimaldot]{babel}
\usepackage{geometry}
\usepackage{array}
\usepackage{tabularx}
\usepackage{xltabular}
\usepackage[table]{xcolor}
\usepackage{fancyhdr}
\usepackage{lastpage}
\usepackage{hyperref}
\usepackage{enumitem}
\usepackage{longtable}
\usepackage{booktabs}
\usepackage{listings}
\usepackage{amssymb}
\usepackage{tikz}
\usetikzlibrary{arrows.meta,shapes.misc,calc}

% ============================================================================
% DISENO VISUAL Y TIPOGRAFIA
% ============================================================================
% Márgenes amplios y proporcionados para lectura técnica en laboratorio.

\geometry{
    left=2cm,
    right=2cm,
    top=3.4cm,
    bottom=2.4cm,
    headheight=58pt,
    headsep=16pt,
    footskip=28pt
}

% Tipografía sans-serif técnica limpia y formal.
\renewcommand{\familydefault}{\sfdefault}

% Paleta institucional de colores sobrios para líneas, tablas y diagramas.
\definecolor{SAIDBlue}{RGB}{35,76,99}
\definecolor{SAIDLight}{RGB}{235,240,243}
\definecolor{SAIDGray}{RGB}{100,100,100}
\definecolor{CodeGray}{RGB}{245,245,245}
\definecolor{SAIDRed}{RGB}{178,34,34}
\definecolor{SAIDGreen}{RGB}{34,139,34}

% Enlaces PDF con color institucional SAIDBlue.
\hypersetup{
    colorlinks=true,
    linkcolor=SAIDBlue,
    urlcolor=SAIDBlue,
    citecolor=SAIDBlue,
    pdftitle={\DocumentTitle\ v\DocumentVersion},
    pdfauthor={\DocumentAuthors},
    pdfsubject={\DocumentSubtitle}
}
\urlstyle{same}

% Formato para fragmentos de código y tramas hexadecimales.
\lstset{
    basicstyle=\ttfamily\small,
    backgroundcolor=\color{CodeGray},
    frame=single,
    columns=fullflexible,
    keepspaces=true,
    showstringspaces=false,
    breaklines=true
}

% Estilos TikZ para diagramas de secuencia claros y consistentes.
\tikzset{
    seqactor/.style={draw=SAIDBlue, fill=SAIDLight, thick, rounded corners=2pt, minimum width=2.8cm, minimum height=0.6cm, font=\bfseries\small},
    seqmsg/.style={-Latex, thick, SAIDBlue, font=\footnotesize},
    seqmsgback/.style={-Latex, thick, SAIDBlue!85!black, font=\footnotesize},
    seqmsglost/.style={-Latex, thick, dashed, SAIDRed, font=\footnotesize},
    seqnote/.style={draw=gray!50, fill=white, thin, font=\scriptsize, align=center, inner sep=2.5pt}
}

% ============================================================================
% COMANDOS / ARTEFACTOS REUTILIZABLES
% ============================================================================
% Comandos normativos reutilizables para garantizar consistencia visual.

% Cabecera de celda para tablas técnicas.
\newcommand{\tablehead}[1]{\cellcolor{SAIDBlue}\color{white}\textbf{#1}}

% Recuadro normativo de requisito verificable obligatorio.
\newcommand{\requirement}[2]{%
    \par\smallskip
    \noindent\begin{tabularx}{\textwidth}{|>{\bfseries}p{3.2cm}|X|}
        \hline
        \cellcolor{SAIDLight}#1 & #2 \\
        \hline
    \end{tabularx}
    \par\smallskip
}

% ============================================================================
% HEADER Y FOOTER
% ============================================================================
% Encabezado y pie de página enlazados directamente a la metadata global.

\pagestyle{fancy}
\fancyhf{}
\renewcommand{\headrulewidth}{0pt}
\renewcommand{\footrulewidth}{0.4pt}

\fancyhead[C]{%
    \begin{tabularx}{\textwidth}{|X|p{2.5cm}|p{2.1cm}|p{2.2cm}|}
        \hline
        \multicolumn{4}{|c|}{\textbf{\DocumentTitle}} \\
        \hline
        \footnotesize \DocumentId &
        \footnotesize Ver. \DocumentVersion &
        \footnotesize \DocumentStatus &
        \footnotesize Pág. \thepage/\pageref*{LastPage} \\
        \hline
    \end{tabularx}%
}

\fancyfoot[L]{\footnotesize \DocumentOwner}
\fancyfoot[R]{\footnotesize \DocumentDate}

% ============================================================================
% CONTENIDO DEL DOCUMENTO
% ============================================================================

\begin{document}

% ----------------------------------------------------------------------------
% PORTADA / FICHA DE CONTROL
% ----------------------------------------------------------------------------
% Ficha formal sobria de identificación técnica del protocolo.

\begin{center}
    {\LARGE\bfseries \DocumentTitle\\[0.3cm]}
    {\large \DocumentSubtitle}
\end{center}

\vspace{0.6cm}

\noindent\begin{tabularx}{\textwidth}{|>{\bfseries}p{4.2cm}|X|}
    \hline
    Identificador & \DocumentId \\
    \hline
    Tipo de documento & \DocumentType \\
    \hline
    Versión & \DocumentVersion \\
    \hline
    Estado & \DocumentStatus \\
    \hline
    Fecha & \DocumentDate \\
    \hline
    Propietario & \DocumentOwner \\
    \hline
    Autor(es) & \DocumentAuthors \\
    \hline
    Consumidores / destinatarios & \DocumentConsumers \\
    \hline
\end{tabularx}

\vspace{0.8cm}

\noindent\textbf{Nota de aplicación técnica.}
Este documento define el formato y la semántica del enlace UART utilizado entre
el microcontrolador ESP32 y el microcontrolador Arduino Mega en el sistema SAID.
La presente especificación en estado \DocumentStatus\ establece los lineamientos
definitivos y cerrados para la implementación de firmware en ambos subsistemas.

\newpage

% ============================================================================
% INDICE DE CONTENIDOS
% ============================================================================

\setcounter{tocdepth}{2}
\begingroup
\makeatletter
\renewcommand*\l@section[2]{%
  \ifnum \c@tocdepth >\z@
    \addpenalty\@secpenalty
    \vskip 0.4em \@plus\p@
    \setlength\@tempdima{2.0em}%
    \begingroup
      \parindent \z@ \rightskip \@pnumwidth
      \parfillskip -\@pnumwidth
      {\leavevmode
       \bfseries #1\nobreak\hfil \nobreak\hb@xt@\@pnumwidth{\hss #2}}\par
    \endgroup
  \fi}
\renewcommand*\l@subsection{\@dottedtocline{2}{2.0em}{2.5em}}
\makeatother
\tableofcontents
\endgroup
\newpage

% ============================================================================
% SECCION 1 - PROPOSITO, ALCANCE Y ARQUITECTURA DEL SISTEMA
% ============================================================================

\section{Propósito, alcance y arquitectura del sistema}

Este documento establece la especificación técnica del contrato lógico de
comunicación serie asíncrona (UART) entre el ESP32 y el Arduino Mega dentro de la
máquina expendedora SAID. Su meta principal es posibilitar la interoperabilidad
directa, confiable y determinista entre el firmware de control de negocio y el
firmware de control físico de periféricos.

\subsection{Distribución de responsabilidades por subsistema}

\noindent\begin{tabularx}{\textwidth}{|p{3.5cm}|X|}
    \hline
    \tablehead{Equipo / Rol} & \tablehead{Responsabilidad técnica en el sistema} \\
    \hline
    Arquitectura & Especifica la estructura de tramas, el catálogo de comandos, las enumeraciones, el orden de bytes y las reglas de idempotencia. \\
    \hline
    Programación ESP32 & Implementa el controlador maestro del protocolo UART, gestiona la máquina de estados de negocio (FSM), consulta y persiste transacciones en SQLite, controla el lector RFID y autoriza las órdenes de venta. \\
    \hline
    Programación Arduino Mega & Implementa el controlador periférico UART, atiende comandos de motor, refresca la pantalla LCD y notifica eventos físicos de teclado matricial y sensores de seguridad. \\
    \hline
    Electrónica & Provee el canal físico serie confiable, selección de líneas TX/RX, niveles de tensión adaptados (3.3~V en ESP32 / 5.0~V en Arduino Mega), plano de masa común y apantallamiento contra ruido de motores. \\
    \hline
\end{tabularx}

\subsection{Arquitectura Web y desacoplamiento de base de datos}

El sistema SAID contempla una interfaz web para tareas de supervisión, recarga
de saldo y mantenimiento administrativo. La relación entre dicha interfaz y el
enlace UART se rige bajo los siguientes principios de arquitectura:

\begin{itemize}[leftmargin=1.6em]
    \item \textbf{Interacción exclusiva vía HTTP con ESP32:} La aplicación web
    se comunica únicamente con el servidor HTTP / API REST que se ejecuta dentro del
    ESP32. En ningún caso la aplicación web interactúa de forma directa con el
    Arduino Mega ni envía tramas UART directas al bus físico.
    \item \textbf{Privacidad de la base de datos SQLite:} La base de datos SQLite
    reside exclusivamente en la memoria de almacenamiento del ESP32. La aplicación web
    accede a los datos mediante endpoints HTTP expuestos por el ESP32, y el Arduino Mega
    no tiene conocimiento alguno de la existencia de tablas relacionales o esquemas de base de datos.
    \item \textbf{Agnosticismo del Arduino Mega ante el origen del comando:}
    El Arduino Mega no distingue si una solicitud de operación (como una prueba de canal
    o una consulta de sensores) fue originada físicamente desde el teclado matricial
    o remotamente desde el panel web. Para el Arduino Mega, toda instrucción es una
    trama UART emitida, validada y autorizada por el ESP32.
\end{itemize}

\subsection{Límites del contrato (Fuera de alcance)}

Los parámetros de la capa física (velocidad en baudios, asignación de pines GPIO,
convertidores de nivel lógico 3.3~V $\leftrightarrow$ 5.0~V y cables) se definen
en la documentación de Electrónica y no forman parte de este contrato lógico.


\clearpage
% ============================================================================
% SECCION 2 - CONVENCIONES DE DATOS
% ============================================================================

\section{Convenciones de datos}

\subsection{Tipos de datos permitidos}

Para prevenir inconsistencias de tamaño entre la arquitectura de 32 bits del
ESP32 y la arquitectura de 8 bits (AVR) del Arduino Mega, todos los campos de
las tramas utilizan exclusivamente tipos binarios de longitud fija de la
biblioteca estándar \texttt{<stdint.h>}.

\noindent\begin{tabularx}{\textwidth}{|p{2.8cm}|p{2.2cm}|X|}
    \hline
    \tablehead{Tipo} & \tablehead{Tamaño} & \tablehead{Uso normativo en el protocolo} \\
    \hline
    \texttt{uint8\_t} & 1 byte (8 bits) & Comandos (\texttt{CMD}), secuencias (\texttt{SEQ}), longitudes (\texttt{LEN}), canales, modos y enumeraciones. \\
    \hline
    \texttt{uint32\_t} & 4 bytes (32 bits) & Identificadores de transacción monotónicos (\texttt{transaction\_id}). \\
    \hline
    ASCII fijo & Longitud fija & Cadenas de caracteres para display LCD sin byte terminador nulo. \\
    \hline
\end{tabularx}

\requirement{UART-REQ-001}{Queda estrictamente prohibido el uso de tipos de datos de ancho indefinido como \texttt{int}, \texttt{short} o \texttt{long} en las estructuras de datos de la interfaz UART.}

\subsection{Orden de bytes (Endianness)}

Todos los campos numéricos enteros de ancho mayor a un byte (\texttt{uint32\_t})
se transmiten en formato \textbf{little-endian} (el byte de menor peso o \textit{Least
Significant Byte} se envía en primer lugar por la línea serie). Por ejemplo, el
identificador de transacción \texttt{0x0000002A} (decimal 42) se transmite como la
secuencia de cuatro bytes: \texttt{0x2A} \texttt{0x00} \texttt{0x00} \texttt{0x00}.

\subsection{Representación de texto}

Los campos de texto utilizados para actualizar la pantalla alfanumérica LCD (\texttt{DISPLAY})
cumplen las siguientes reglas:
\begin{itemize}[leftmargin=1.6em,itemsep=1pt,parsep=0pt,topsep=2pt]
    \item Codificación estándar en ASCII imprimible de 7 bits (\texttt{0x20} a \texttt{0x7E}).
    \item Cada línea de la pantalla tiene una longitud exacta de 16 caracteres.
    \item No se utiliza carácter terminador nulo (\lstinline|'\0'|).
    \item Cadenas con menos de 16 caracteres se rellenan a la derecha con espacios ASCII (\texttt{0x20}).
    \item Caracteres especiales o acentuados fuera del rango ASCII imprimible deben sustituirse por el carácter \texttt{'?'} (\texttt{0x3F}).
\end{itemize}

\subsection{Identificador de transacción (\texttt{transaction\_id})}

El identificador \texttt{transaction\_id} de 32 bits vincula unívocamente las órdenes
de dispensado físico con su registro contable:
\begin{itemize}[leftmargin=1.6em,itemsep=1pt,parsep=0pt,topsep=2pt]
    \item \textbf{Generador único:} El ESP32 es el único dispositivo facultado para generar nuevos \texttt{transaction\_id}.
    \item \textbf{Rango operativo:} Valores válidos comprendidos entre \texttt{0x00000001} y \texttt{0xFFFFFFFE}.
    \item \textbf{Valores centinela:} \texttt{0x00000000} (inválida o nula) y \texttt{0xFFFFFFFF} (sin transacción en curso).
\end{itemize}

\subsection{Número de secuencia (\texttt{SEQ})}

Cada trama transmitida incluye un campo \texttt{SEQ} de 8 bits (\texttt{uint8\_t}):
\begin{itemize}[leftmargin=1.6em,itemsep=1pt,parsep=0pt,topsep=2pt]
    \item Se incrementa progresivamente módulo 256 (\texttt{0} $\rightarrow$ \texttt{1} $\rightarrow$ \dots $\rightarrow$ \texttt{255} $\rightarrow$ \texttt{0}) con cada nuevo mensaje generado.
    \item Las respuestas directas e inmediatas (\texttt{ACK}, respuesta de \texttt{STATUS}, eco de \texttt{HEARTBEAT} y respuesta de \texttt{HELLO}) replican con exactitud el valor de \texttt{SEQ} del mensaje entrante que las provocó.
\end{itemize}

\clearpage


% ============================================================================
% SECCION 3 - FORMATO GENERAL DEL FRAME
% ============================================================================

\section{Formato general del frame}

Todas las tramas transmitidas entre el ESP32 y el Arduino Mega comparten una
estructura binaria regular y uniforme delimitada por un preámbulo fijo de dos bytes
y un campo de longitud explícito.

\subsection{Estructura binaria de la trama}

\begin{center}
\begin{tabular}{|c|c|c|c|c|c|c|}
    \hline
    \textbf{SOF1} & \textbf{SOF2} & \textbf{PROTO} & \textbf{CMD} & \textbf{SEQ} & \textbf{LEN} & \textbf{PAYLOAD} \\
    \hline
    1 byte & 1 byte & 1 byte & 1 byte & 1 byte & 1 byte & 0 a 32 bytes \\
    \hline
\end{tabular}
\end{center}

\subsection{Campos del frame}

\noindent\begin{xltabular}{\textwidth}{|p{1.8cm}|p{1.5cm}|p{2.0cm}|X|}
    \hline
    \tablehead{Campo} & \tablehead{Offset} & \tablehead{Tamaño} & \tablehead{Descripción y reglas} \\
    \hline \endhead

    \texttt{SOF1} & 0 & 1 byte & Primer byte de sincronización. Constante fija \texttt{0xA5}. \\
    \hline
    \texttt{SOF2} & 1 & 1 byte & Segundo byte de sincronización. Constante fija \texttt{0x5A}. \\
    \hline
    \texttt{PROTO} & 2 & 1 byte & Versión mayor del protocolo a nivel de trama. Constante \texttt{0x02}. \\
    \hline
    \texttt{CMD} & 3 & 1 byte & Identificador del comando según el catálogo oficial. \\
    \hline
    \texttt{SEQ} & 4 & 1 byte & Número de secuencia del mensaje ($0 \le \texttt{SEQ} \le 255$). \\
    \hline
    \texttt{LEN} & 5 & 1 byte & Longitud en bytes del campo \texttt{PAYLOAD} ($0 \le \texttt{LEN} \le 32$). \\
    \hline
    \texttt{PAYLOAD} & 6 & 0..32 bytes & Datos asociados al comando. Si $\texttt{LEN} = 0$, este campo no existe. \\
    \hline
\end{xltabular}

\subsection{Reglas de validación y delimitación}

\begin{itemize}[leftmargin=1.6em]
    \item \textbf{Límites dimensionales:} La trama mínima sin carga útil ocupa exactamente 6 bytes ($\texttt{LEN} = 0$). La trama de máxima longitud ocupa 38 bytes ($\texttt{LEN} = 32$).
    \item \textbf{Rechazo por desbordamiento de longitud:} Si el receptor procesa un byte \texttt{LEN} estrictamente mayor a 32 ($\texttt{LEN} > 32$), la trama se descarta inmediatamente por considerarse corrupta o desfasada.
    \item \textbf{Ausencia de checksum / CRC:} Esta versión 2.0.0 del contrato no utiliza código de redundancia cíclica (CRC) ni checksum adicional; la delimitación y validación se basan exclusivamente en el preámbulo fijo de dos bytes, el control de versión de protocolo y el campo de longitud explícito.
    \item \textbf{Resincronización del parser:} Ante la recepción de bytes corruptos, discrepancia en \texttt{PROTO} ($\neq \texttt{0x02}$) o valor inválido de \texttt{LEN}, el receptor limpia su buffer temporal de recepción y reanuda la búsqueda de la secuencia de sincronización \texttt{0xA5 0x5A}.
\end{itemize}


% ============================================================================
% SECCION 4 - DICCIONARIO DEFINITIVO DE COMANDOS
% ============================================================================

\clearpage
\section{Diccionario definitivo de comandos}

El protocolo SAID v2.0.0 define un conjunto cerrado de 9 comandos que cubren la
totalidad de los intercambios requeridos para el control de la máquina expendedora:

\noindent\begin{xltabular}{\textwidth}{|p{1.4cm}|p{2.5cm}|p{3.2cm}|X|}
    \hline
    \tablehead{CMD} & \tablehead{Nombre} & \tablehead{Dirección} & \tablehead{Semántica operativa} \\
    \hline \endhead

    \texttt{0x01} & \texttt{HELLO} & Ambas direcciones & Sincronización inicial de arranque y verificación de versión de protocolo. \\
    \hline
    \texttt{0x02} & \texttt{ACK} & Ambas direcciones & Confirmación inmediata de recepción válida o notificación de rechazo con motivo. \\
    \hline
    \texttt{0x03} & \texttt{STATUS} & Ambas direcciones & Solicitud (ESP32 $\rightarrow$ Mega) y reporte (Mega $\rightarrow$ ESP32) de sensores y modo. \\
    \hline
    \texttt{0x04} & \texttt{SET\_MODE} & ESP32 $\rightarrow$ Mega & Configuración del modo operativo autorizado en Arduino Mega (Venta o Mantenimiento). \\
    \hline
    \texttt{0x05} & \texttt{HEARTBEAT} & ESP32 $\rightarrow$ Mega $\rightarrow$ ESP32 & Diagnóstico puntual del canal serie mediante patrón simple de ping/eco. \\
    \hline
    \texttt{0x10} & \texttt{KEY} & Mega $\rightarrow$ ESP32 & Notificación de evento de tecla presionada tras filtrado debounce en hardware local. \\
    \hline
    \texttt{0x11} & \texttt{DISPLAY} & ESP32 $\rightarrow$ Mega & Actualización del contenido textual completo (2 líneas de 16 caracteres) en el display LCD. \\
    \hline
    \texttt{0x12} & \texttt{VEND} & ESP32 $\rightarrow$ Mega & Orden de dispensado físico autorizada para un canal mecánico determinado. \\
    \hline
    \texttt{0x13} & \texttt{RESULT} & Mega $\rightarrow$ ESP32 & Notificación asíncrona del resultado físico del dispensado ejecutado en Mega. \\
    \hline
\end{xltabular}

\vspace{0.4cm}
\noindent\textbf{Cláusula de extensión futura:}
Los códigos de comando no listados en la tabla anterior quedan reservados para futuras
versiones del protocolo.

% ============================================================================
% SECCION 5 - ENUMERACIONES COMPARTIDAS
% ============================================================================

\section{Enumeraciones compartidas}

\subsection{Resultado de confirmación (\texttt{resultado} en comando ACK)}

\noindent\begin{tabularx}{\textwidth}{|p{1.8cm}|p{3.5cm}|X|}
    \hline
    \tablehead{Valor} & \tablehead{Constante} & \tablehead{Definición} \\
    \hline
    \texttt{0x00} & \texttt{ACK\_RECEIVED} & Trama recibida íntegra, sintácticamente válida y aceptada para ejecución. \\
    \hline
    \texttt{0x01} & \texttt{ACK\_REJECTED} & Trama rechazada por incoherencia sintáctica, estado inválido o conflicto local. \\
    \hline
\end{tabularx}

\requirement{UART-REQ-002}{\texttt{ACK\_RECEIVED} certifica la recepción y aceptación lógica del mensaje en el receptor; no constituye una confirmación de que la acción física mecánica haya concluido.}

\subsection{Motivos de rechazo (\texttt{motivo} en comando ACK)}

\noindent\begin{xltabular}{\textwidth}{|p{1.8cm}|p{4.8cm}|X|}
    \hline
    \tablehead{Valor} & \tablehead{Constante} & \tablehead{Causa reportada} \\
    \hline \endhead

    \texttt{0x00} & \texttt{NONE} & Sin error. Empleado cuando \texttt{resultado} es \texttt{ACK\_RECEIVED}. \\
    \hline
    \texttt{0x01} & \texttt{INVALID\_CMD} & El comando recibido no existe en el catálogo oficial. \\
    \hline
    \texttt{0x02} & \texttt{INVALID\_LENGTH} & El campo \texttt{LEN} recibido no coincide con el tamaño esperado para el comando. \\
    \hline
    \texttt{0x03} & \texttt{INVALID\_PAYLOAD} & La carga útil contiene parámetros fuera de rango válido. \\
    \hline
    \texttt{0x04} & \texttt{INVALID\_STATE} & El comando no está permitido en el estado actual de operación. \\
    \hline
    \texttt{0x05} & \texttt{BUSY} & El subsistema se encuentra ocupado ejecutando una operación previa. \\
    \hline
    \texttt{0x06} & \texttt{INVALID\_CHANNEL} & El canal de motor solicitado no pertenece al rango admitido (1 a 4). \\
    \hline
    \texttt{0x07} & \texttt{VERSION\_MISMATCH} & Discrepancia incompatible en versión de protocolo negociada en \texttt{HELLO}. \\
    \hline
    \texttt{0x08} & \texttt{DUPLICATE\_CONFLICT} & Se recibió una solicitud con idéntico identificador pero parámetros contradictorios. \\
    \hline
    \texttt{0x09} & \texttt{UNSAFE\_PHYSICAL\_STATE} & Condición física que impide operar con seguridad (p. ej., puerta abierta). \\
    \hline
    \texttt{0x0A} & \texttt{INTERNAL\_ERROR} & Falla interna de hardware o periférico en el nodo receptor. \\
    \hline
    \texttt{0x0B} & \texttt{UNKNOWN\_TRANSACTION} & Transacción referenciada no reconocida en el nodo. \\
    \hline
\end{xltabular}

\subsection{Resultado físico de dispensado (\texttt{result} en comando RESULT)}

\noindent\begin{tabularx}{\textwidth}{|p{1.8cm}|p{5.0cm}|X|}
    \hline
    \tablehead{Valor} & \tablehead{Constante} & \tablehead{Definición física observable} \\
    \hline
    \texttt{0x00} & \texttt{DELIVERED} & El motor completó el ciclo y el sensor de caída detectó la entrega del producto. \\
    \hline
    \texttt{0x01} & \texttt{REJECTED\_BEFORE\_MOTION} & La orden fue rechazada antes de activar actuadores o motores mecánicos. \\
    \hline
    \texttt{0x02} & \texttt{UNCERTAIN} & El resultado físico no pudo determinarse de forma confiable (p. ej., bloqueo mecánico o falla de sensor). \\
    \hline
\end{tabularx}

\subsection{Identificador de rol (\texttt{role} en comando HELLO)}

\noindent\begin{tabularx}{\textwidth}{|p{1.8cm}|p{3.2cm}|X|}
    \hline
    \tablehead{Valor} & \tablehead{Constante} & \tablehead{Rol asignado en el sistema} \\
    \hline
    \texttt{0x00} & \texttt{ROLE\_ESP32} & Nodo maestro (controlador de lógica de negocio, SQLite y RFID). \\
    \hline
    \texttt{0x01} & \texttt{ROLE\_MEGA} & Nodo periférico (controlador de teclado, display LCD, motores y sensores). \\
    \hline
\end{tabularx}

\subsection{Modo de operación del sistema (\texttt{mode})}

\noindent\begin{tabularx}{\textwidth}{|p{1.8cm}|p{4.5cm}|X|}
    \hline
    \tablehead{Valor} & \tablehead{Constante} & \tablehead{Comportamiento operativo} \\
    \hline
    \texttt{0x00} & \texttt{MODE\_VENTA} & Modo regular de atención al público: teclado habilitado y órdenes \texttt{VEND} admitidas. \\
    \hline
    \texttt{0x01} & \texttt{MODE\_MANTENIMIENTO} & Modo técnico de servicio y diagnóstico: venta al público deshabilitada. \\
    \hline
\end{tabularx}

\requirement{UART-REQ-003}{Al energizarse o experimentar un reinicio, el Arduino Mega inicia de forma obligatoria en \texttt{MODE\_MANTENIMIENTO} como medida de seguridad física, permaneciendo en dicho modo hasta que el ESP32 emita explícitamente el comando \texttt{SET\_MODE(MODE\_VENTA)}.}

\subsection{Estados de sensores físicos}

Valores discretos reportados por el Arduino Mega en la respuesta del comando \texttt{STATUS}:
\begin{itemize}[leftmargin=1.6em]
    \item \textbf{Sensor de puerta (\texttt{door}):} \texttt{0x00} = \texttt{CLOSED} (puerta cerrada y segura), \texttt{0x01} = \texttt{OPEN} (puerta abierta), \texttt{0xFF} = \texttt{UNKNOWN}.
    \item \textbf{Barrera óptica de caída (\texttt{barrier}):} \texttt{0x00} = \texttt{FREE} (haz despejado), \texttt{0x01} = \texttt{OCCUPIED} (haz obstruido), \texttt{0xFF} = \texttt{UNKNOWN}.
\end{itemize}


% ============================================================================
% SECCION 6 - ESPECIFICACION DETALLADA DE MENSAJES
% ============================================================================

\section{Especificación de mensajes}

\subsection{\texttt{HELLO} -- Inicialización y sincronización de versión (\texttt{0x01})}

\noindent
\textbf{Dirección:} Ambas direcciones (ESP32 $\leftrightarrow$ Arduino Mega).\\
\textbf{Longitud de payload:} Exactamente 3 bytes ($\texttt{LEN} = 3$).\\
\textbf{Propósito:} Handshake de inicio para comprobar compatibilidad mutua de versión y roles.

\noindent\begin{xltabular}{\textwidth}{|p{1.1cm}|p{1.6cm}|p{1.7cm}|p{4.2cm}|X|}
    \hline
    \tablehead{Offset} & \tablehead{Tamaño} & \tablehead{Tipo} & \tablehead{Campo} & \tablehead{Descripción} \\
    \hline \endhead

    0 & 1 byte & \texttt{uint8\_t} & \texttt{version\_major} & Versión mayor del protocolo. Valor constante \texttt{2}. \\
    \hline
    1 & 1 byte & \texttt{uint8\_t} & \texttt{version\_minor} & Versión menor del protocolo. Valor constante \texttt{0}. \\
    \hline
    2 & 1 byte & \texttt{uint8\_t} & \texttt{role} & Identificador del nodo emisor (\texttt{0x00} ESP32, \texttt{0x01} Mega). \\
    \hline
\end{xltabular}

\requirement{UART-REQ-004}{Todo mensaje \texttt{HELLO} recibido debe responderse de inmediato con otro mensaje \texttt{HELLO} emitido por el receptor, conservando exactamente el mismo número de secuencia \texttt{SEQ}.}

\subsection{\texttt{ACK} -- Confirmación inmediata (\texttt{0x02})}

\noindent
\textbf{Dirección:} Ambas direcciones (ESP32 $\leftrightarrow$ Arduino Mega).\\
\textbf{Longitud de payload:} Exactamente 3 bytes ($\texttt{LEN} = 3$).\\
\textbf{Propósito:} Acusar recibo inmediato de una solicitud entrante indicando si fue aceptada o rechazada.

\noindent\begin{xltabular}{\textwidth}{|p{1.1cm}|p{1.6cm}|p{1.7cm}|p{4.2cm}|X|}
    \hline
    \tablehead{Offset} & \tablehead{Tamaño} & \tablehead{Tipo} & \tablehead{Campo} & \tablehead{Descripción} \\
    \hline \endhead

    0 & 1 byte & \texttt{uint8\_t} & \texttt{cmd\_referenciado} & Identificador del comando al cual responde. \\
    \hline
    1 & 1 byte & \texttt{uint8\_t} & \texttt{resultado} & \texttt{0x00} (\texttt{ACK\_RECEIVED}) o \texttt{0x01} (\texttt{ACK\_REJECTED}). \\
    \hline
    2 & 1 byte & \texttt{uint8\_t} & \texttt{motivo} & Motivo normativo de rechazo (\texttt{0x00} si fue aceptado). \\
    \hline
\end{xltabular}

\requirement{UART-REQ-005}{Una trama de confirmación \texttt{ACK} nunca genera otro \texttt{ACK} en respuesta, evitando la formación de bucles infinitos en el canal serie.}

\subsection{\texttt{STATUS} -- Consulta y reporte de estado (\texttt{0x03})}

\noindent
\textbf{Dirección:} Solicitud: ESP32 $\rightarrow$ Mega; Respuesta: Mega $\rightarrow$ ESP32.\\
\textbf{Propósito:} Intercambiar el estado de sensores físicos y el modo activo del Arduino Mega.
\begin{itemize}[leftmargin=1.6em]
    \item \textbf{Solicitud de estado:} Trama emitida por ESP32 con payload vacío ($\texttt{LEN} = 0$).
    \item \textbf{Respuesta de estado:} Trama emitida por Mega con payload de 7 bytes ($\texttt{LEN} = 7$) y el mismo \texttt{SEQ} de la solicitud.
\end{itemize}

\paragraph{Payload de la respuesta \texttt{STATUS} (7 bytes):}

\noindent\begin{xltabular}{\textwidth}{|p{1.1cm}|p{1.6cm}|p{1.7cm}|p{4.5cm}|X|}
    \hline
    \tablehead{Offset} & \tablehead{Tamaño} & \tablehead{Tipo} & \tablehead{Campo} & \tablehead{Descripción} \\
    \hline \endhead

    0 & 1 byte & \texttt{uint8\_t} & \texttt{door} & Estado del sensor de puerta (\texttt{0x00} cerrada, \texttt{0x01} abierta). \\
    \hline
    1 & 1 byte & \texttt{uint8\_t} & \texttt{barrier} & Estado de la barrera óptica (\texttt{0x00} libre, \texttt{0x01} ocupada). \\
    \hline
    2 & 1 byte & \texttt{uint8\_t} & \texttt{mode} & Modo operativo actual en Mega (\texttt{0x00} Venta, \texttt{0x01} Mantenimiento). \\
    \hline
    3 & 4 bytes & \texttt{uint32\_t} & \texttt{current\_transaction\_id} & Transacción en ejecución (\texttt{0xFFFFFFFF} si inactivo). \\
    \hline
\end{xltabular}

\requirement{UART-REQ-006}{El comando \texttt{STATUS} no requiere tramas \texttt{ACK}. La propia trama de respuesta \texttt{STATUS} ($\texttt{LEN} = 7$) sirve como confirmación de la consulta.}

\subsection{\texttt{SET\_MODE} -- Configuración de modo operativo (\texttt{0x04})}

\noindent
\textbf{Dirección:} ESP32 $\rightarrow$ Arduino Mega.\\
\textbf{Longitud de payload:} Exactamente 1 byte ($\texttt{LEN} = 1$).\\
\textbf{Propósito:} Comandar a Arduino Mega la transición entre modo de venta y modo de mantenimiento.

\noindent\begin{xltabular}{\textwidth}{|p{1.1cm}|p{1.6cm}|p{1.7cm}|p{4.2cm}|X|}
    \hline
    \tablehead{Offset} & \tablehead{Tamaño} & \tablehead{Tipo} & \tablehead{Campo} & \tablehead{Descripción} \\
    \hline \endhead

    0 & 1 byte & \texttt{uint8\_t} & \texttt{mode} & Modo solicitado: \texttt{0x00} (\texttt{MODE\_VENTA}) o \texttt{0x01} (\texttt{MODE\_MANTENIMIENTO}). \\
    \hline
\end{xltabular}

\requirement{UART-REQ-007}{Si el Arduino Mega se encuentra ejecutando físicamente un dispensado mecánico cuando recibe \texttt{SET\_MODE}, debe responder con \texttt{ACK\_REJECTED} y motivo \texttt{BUSY}, reteniendo el modo activo sin interrumpir el movimiento en curso.}

\subsection{\texttt{HEARTBEAT} -- Verificación de enlace (\texttt{0x05})}

\noindent
\textbf{Dirección:} ESP32 $\rightarrow$ Arduino Mega $\rightarrow$ ESP32.\\
\textbf{Longitud de payload:} 0 bytes ($\texttt{LEN} = 0$).\\
\textbf{Propósito:} Mensaje de diagnóstico simple tipo ping/eco originado por el ESP32 para auditar la conectividad básica del enlace. El Arduino Mega responde emitiendo una trama idéntica en modo eco utilizando el mismo valor de \texttt{SEQ}. No existen reglas de supervisión periódica obligatoria ni estados de enlace degradado definidos normativamente.

\subsection{\texttt{KEY} -- Evento de teclado matricial (\texttt{0x10})}

\noindent
\textbf{Dirección:} Arduino Mega $\rightarrow$ ESP32.\\
\textbf{Longitud de payload:} Exactamente 2 bytes ($\texttt{LEN} = 2$).\\
\textbf{Propósito:} Informar la pulsación estable de una tecla tras filtrado por debounce. El Arduino Mega no asigna significado de negocio al dígito; el ESP32 interpreta su efecto en la máquina de estados.

\noindent\begin{xltabular}{\textwidth}{|p{1.1cm}|p{1.6cm}|p{1.7cm}|p{4.2cm}|X|}
    \hline
    \tablehead{Offset} & \tablehead{Tamaño} & \tablehead{Tipo} & \tablehead{Campo} & \tablehead{Descripción} \\
    \hline \endhead

    0 & 1 byte & \texttt{uint8\_t} & \texttt{key\_ascii} & Código ASCII de la tecla pulsada (\texttt{'0'}..\texttt{'9'}, \texttt{'A'}..\texttt{'D'}, \texttt{'*'}, \texttt{'\#'}). \\
    \hline
    1 & 1 byte & \texttt{uint8\_t} & \texttt{key\_seq} & Contador secuencial de tecla local (módulo 256) generado por Mega. \\
    \hline
\end{xltabular}

\requirement{UART-REQ-008}{El ESP32 debe responder de inmediato a cada trama \texttt{KEY} recibida con un mensaje \texttt{ACK}.}

\subsection{\texttt{DISPLAY} -- Actualización de pantalla LCD (\texttt{0x11})}

\noindent
\textbf{Dirección:} ESP32 $\rightarrow$ Arduino Mega.\\
\textbf{Longitud de payload:} Exactamente 32 bytes ($\texttt{LEN} = 32$).\\
\textbf{Propósito:} Actualizar el contenido textual de las 2 líneas de 16 caracteres del LCD.

\noindent\begin{xltabular}{\textwidth}{|p{1.1cm}|p{1.6cm}|p{1.7cm}|p{4.2cm}|X|}
    \hline
    \tablehead{Offset} & \tablehead{Tamaño} & \tablehead{Tipo} & \tablehead{Campo} & \tablehead{Descripción} \\
    \hline \endhead

    0 & 16 bytes & ASCII fijo & \texttt{line1[16]} & Texto de la línea superior (16 caracteres, relleno con espacios). \\
    \hline
    16 & 16 bytes & ASCII fijo & \texttt{line2[16]} & Texto de la línea inferior (16 caracteres, relleno con espacios). \\
    \hline
\end{xltabular}

\requirement{UART-REQ-009}{Arduino Mega no mantiene menús propios ni interpreta el contenido del display: copia directamente los 32 caracteres en el controlador del LCD y responde de inmediato con \texttt{ACK\_RECEIVED}.}

\subsection{\texttt{VEND} -- Orden de dispensado físico (\texttt{0x12})}

\noindent
\textbf{Dirección:} ESP32 $\rightarrow$ Arduino Mega.\\
\textbf{Longitud de payload:} Exactamente 5 bytes ($\texttt{LEN} = 5$).\\
\textbf{Propósito:} Comandar el accionamiento del motor de entrega de un canal físico asociado a una transacción debidamente autorizada.

\noindent\begin{xltabular}{\textwidth}{|p{1.1cm}|p{1.6cm}|p{1.7cm}|p{4.2cm}|X|}
    \hline
    \tablehead{Offset} & \tablehead{Tamaño} & \tablehead{Tipo} & \tablehead{Campo} & \tablehead{Descripción} \\
    \hline \endhead

    0 & 1 byte & \texttt{uint8\_t} & \texttt{channel} & Canal físico del producto a dispensar (valores válidos: 1 a 4). \\
    \hline
    1 & 4 bytes & \texttt{uint32\_t} & \texttt{transaction\_id} & Identificador único de transacción generado por el ESP32. \\
    \hline
\end{xltabular}

\paragraph{Validaciones previas en Arduino Mega:}
Antes de activar el actuador y responder con \texttt{ACK\_RECEIVED}, el Arduino Mega valida:
\begin{enumerate}[leftmargin=1.8em]
    \item Que el modo del sistema sea estrictamente \texttt{MODE\_VENTA}.
    \item Que el canal recibido pertenezca al rango de hardware conectado ($1 \le \texttt{channel} \le 4$).
    \item Que ningún motor se encuentre en movimiento activo en ese instante.
    \item Que el sensor de puerta confirme que la puerta se encuentra cerrada.
    \item Que el \texttt{transaction\_id} no sea nulo (\texttt{0x00000000}) ni centinela (\texttt{0xFFFFFFFF}).
\end{enumerate}

\requirement{UART-REQ-010}{Regla de idempotencia fundamental: Un mismo identificador \texttt{transaction\_id} jamás debe provocar el accionamiento del motor por segunda vez. Si el Arduino Mega recibe un comando \texttt{VEND} con un \texttt{transaction\_id} idéntico al de una orden ya procesada o en curso, no debe reiniciar ni repetir el movimiento mecánico; responderá de inmediato con \texttt{ACK\_RECEIVED} y reenviará el resultado físico obtenido previamente.}

\requirement{UART-REQ-011}{Si arriba un comando \texttt{VEND} con un \texttt{transaction\_id} idéntico a uno previo pero indicando un canal distinto, el Arduino Mega rechazará la trama inmediatamente con \texttt{ACK\_REJECTED} y motivo \texttt{DUPLICATE\_CONFLICT}.}

\subsection{\texttt{RESULT} -- Reporte de resultado físico (\texttt{0x13})}

\noindent
\textbf{Dirección:} Arduino Mega $\rightarrow$ ESP32.\\
\textbf{Longitud de payload:} Exactamente 5 bytes ($\texttt{LEN} = 5$).\\
\textbf{Propósito:} Notificación asíncrona enviada por el Arduino Mega una vez concluido el ciclo de dispensado para reportar el resultado de la entrega física.

\noindent\begin{xltabular}{\textwidth}{|p{1.1cm}|p{1.6cm}|p{1.7cm}|p{4.2cm}|X|}
    \hline
    \tablehead{Offset} & \tablehead{Tamaño} & \tablehead{Tipo} & \tablehead{Campo} & \tablehead{Descripción} \\
    \hline \endhead

    0 & 4 bytes & \texttt{uint32\_t} & \texttt{transaction\_id} & Identificador de la transacción que motivó la entrega. \\
    \hline
    4 & 1 byte & \texttt{uint8\_t} & \texttt{result} & Desenlace físico: \texttt{DELIVERED}, \texttt{REJECTED\_BEFORE\_MOTION} o \texttt{UNCERTAIN}. \\
    \hline
\end{xltabular}

\requirement{UART-REQ-012}{El ESP32 debe confirmar la recepción del mensaje \texttt{RESULT} emitiendo de inmediato un mensaje \texttt{ACK}.}


% ============================================================================
% SECCION 7 - MATRIZ NORMATIVA DE RESPUESTAS
% ============================================================================

\clearpage
\section{Matriz de respuestas del protocolo}

La siguiente matriz resume la respuesta inmediata obligatoria ante la recepción
de cada comando del protocolo:

\noindent\begin{xltabular}{\textwidth}{|p{3.2cm}|p{3.5cm}|X|}
    \hline
    \tablehead{Mensaje recibido} & \tablehead{Dirección del mensaje} & \tablehead{Respuesta inmediata obligatoria} \\
    \hline \endhead

    \texttt{HELLO} & Ambas direcciones & \texttt{HELLO} (mismo \texttt{SEQ}, confirma versión y rol). \\
    \hline
    \texttt{ACK} & Ambas direcciones & \textbf{Ninguna} (nunca se responde a un ACK). \\
    \hline
    \texttt{STATUS} (solicitud) & ESP32 $\rightarrow$ Mega & \texttt{STATUS} (respuesta con $\texttt{LEN} = 7$, mismo \texttt{SEQ}). \\
    \hline
    \texttt{SET\_MODE} & ESP32 $\rightarrow$ Mega & \texttt{ACK} (\texttt{RECEIVED} o \texttt{REJECTED} con motivo). \\
    \hline
    \texttt{HEARTBEAT} & ESP32 $\rightarrow$ Mega & \texttt{HEARTBEAT} (modo eco, $\texttt{LEN} = 0$, mismo \texttt{SEQ}). \\
    \hline
    \texttt{KEY} & Mega $\rightarrow$ ESP32 & \texttt{ACK} (\texttt{RECEIVED} o \texttt{REJECTED}). \\
    \hline
    \texttt{DISPLAY} & ESP32 $\rightarrow$ Mega & \texttt{ACK} (\texttt{RECEIVED} o \texttt{REJECTED}). \\
    \hline
    \texttt{VEND} & ESP32 $\rightarrow$ Mega & \texttt{ACK} inmediato + \texttt{RESULT} asíncrono posterior. \\
    \hline
    \texttt{RESULT} & Mega $\rightarrow$ ESP32 & \texttt{ACK} (\texttt{RECEIVED} o \texttt{REJECTED}). \\
    \hline
\end{xltabular}


% ============================================================================
% SECCION 8 - DIAGRAMAS DE SECUENCIA DE REFERENCIA
% ============================================================================

\section{Diagramas de secuencia de referencia}

Los siguientes diagramas ilustran los flujos operativos estándar entre el ESP32
y el Arduino Mega para los escenarios cardinales del sistema.

\subsection{Secuencia 1: Arranque normal y configuración de modo}

\begin{center}
\begin{tikzpicture}
    \node[seqactor] at (2.8,0) (esp) {ESP32};
    \node[seqactor] at (9.8,0) (mega) {Arduino Mega};
    \draw[thick, dashed, gray!60] (2.8,-0.3) -- (2.8,-3.8);
    \draw[thick, dashed, gray!60] (9.8,-0.3) -- (9.8,-3.8);

    \node[seqnote, anchor=west] at (10.0,-0.6) {Mega inicia por defecto\\en MODE\_MANTENIMIENTO};
    \draw[seqmsg] (2.8,-1.2) -- node[above] {HELLO [SEQ=0x01, v2.0, ESP32]} (9.8,-1.2);
    \draw[seqmsgback] (9.8,-1.8) -- node[above] {HELLO [SEQ=0x01, v2.0, MEGA]} (2.8,-1.8);
    \node[seqnote] at (6.3,-2.3) {Handshake de versiones 2.0 completado con éxito};
    \draw[seqmsg] (2.8,-2.9) -- node[above] {SET\_MODE [SEQ=0x02, MODE\_VENTA]} (9.8,-2.9);
    \draw[seqmsgback] (9.8,-3.5) -- node[above] {ACK [SEQ=0x02, cmd=0x04, RECEIVED, NONE]} (2.8,-3.5);
\end{tikzpicture}
\end{center}

\subsection{Secuencia 2: Captura y notificación de tecla física}

\begin{center}
\begin{tikzpicture}
    \node[seqactor] at (2.8,0) (esp) {ESP32};
    \node[seqactor] at (9.8,0) (mega) {Arduino Mega};
    \draw[thick, dashed, gray!60] (2.8,-0.3) -- (2.8,-2.6);
    \draw[thick, dashed, gray!60] (9.8,-0.3) -- (9.8,-2.6);

    \node[seqnote, anchor=west] at (10.0,-0.7) {Usuario presiona tecla `5'\\Filtro debounce validado};
    \draw[seqmsgback] (9.8,-1.3) -- node[above] {KEY [SEQ=0x14, ascii=`5' (0x35), key\_seq=0x07]} (2.8,-1.3);
    \draw[seqmsg] (2.8,-1.9) -- node[above] {ACK [SEQ=0x14, cmd=0x10, RECEIVED, NONE]} (9.8,-1.9);
    \node[seqnote, anchor=east] at (2.6,-2.3) {ESP32 evalúa contexto en FSM};
\end{tikzpicture}
\end{center}

\subsection{Secuencia 3: Venta exitosa completa}

\begin{center}
\begin{tikzpicture}
    \node[seqactor] at (2.8,0) (esp) {ESP32};
    \node[seqactor] at (9.8,0) (mega) {Arduino Mega};
    \draw[thick, dashed, gray!60] (2.8,-0.3) -- (2.8,-4.2);
    \draw[thick, dashed, gray!60] (9.8,-0.3) -- (9.8,-4.2);

    \node[seqnote, anchor=east] at (2.6,-0.6) {ESP32 autoriza compra\\y genera tx\_id=42};
    \draw[seqmsg] (2.8,-1.0) -- node[above] {VEND [SEQ=0x20, ch=3, tx=42]} (9.8,-1.0);
    \draw[seqmsgback] (9.8,-1.6) -- node[above] {ACK [SEQ=0x20, cmd=0x12, RECEIVED, NONE]} (2.8,-1.6);
    \node[seqnote, anchor=west] at (10.0,-2.1) {Mega acciona motor canal 3;\\sensor óptico detecta caída};
    \draw[seqmsgback] (9.8,-2.8) -- node[above] {RESULT [SEQ=0x21, tx=42, DELIVERED]} (2.8,-2.8);
    \draw[seqmsg] (2.8,-3.4) -- node[above] {ACK [SEQ=0x21, cmd=0x13, RECEIVED, NONE]} (9.8,-3.4);
    \node[seqnote, anchor=east] at (2.6,-3.9) {ESP32 asienta venta en SQLite};
\end{tikzpicture}
\end{center}


\subsection{Secuencia 4: Idempotencia ante VEND duplicado}

\begin{center}
\begin{tikzpicture}
    \node[seqactor] at (2.8,0) (esp) {ESP32};
    \node[seqactor] at (9.8,0) (mega) {Arduino Mega};
    \draw[thick, dashed, gray!60] (2.8,-0.3) -- (2.8,-4.2);
    \draw[thick, dashed, gray!60] (9.8,-0.3) -- (9.8,-4.2);

    \draw[seqmsg] (2.8,-0.8) -- node[above] {VEND [SEQ=0x20, ch=3, tx=42]} (9.8,-0.8);
    \draw[seqmsglost] (9.8,-1.4) -- (6.3,-1.4) node[midway, above] {ACK inicial perdido} node[SAIDRed, font=\bfseries] {$\times$};
    \node[seqnote, anchor=east] at (2.6,-2.0) {ESP32 retransmite solicitud\\con idéntico SEQ y tx\_id=42};
    \draw[seqmsg] (2.8,-2.6) -- node[above] {VEND [SEQ=0x20, ch=3, tx=42]} (9.8,-2.6);
    \node[seqnote, anchor=west] at (10.0,-2.9) {Mega detecta tx\_id=42 repetida:\\\textbf{NO acciona el motor por 2da vez}};
    \draw[seqmsgback] (9.8,-3.6) -- node[above] {ACK [SEQ=0x20, cmd=0x12, RECEIVED, NONE]} (2.8,-3.6);
\end{tikzpicture}
\end{center}

\subsection{Secuencia 5: Consulta de estado de sensores (STATUS)}

\begin{center}
\begin{tikzpicture}
    \node[seqactor] at (2.8,0) (esp) {ESP32};
    \node[seqactor] at (9.8,0) (mega) {Arduino Mega};
    \draw[thick, dashed, gray!60] (2.8,-0.3) -- (2.8,-2.6);
    \draw[thick, dashed, gray!60] (9.8,-0.3) -- (9.8,-2.6);

    \draw[seqmsg] (2.8,-0.8) -- node[above] {STATUS [SEQ=0x25, LEN=0]} (9.8,-0.8);
    \draw[seqmsgback] (9.8,-1.6) -- node[above] {STATUS resp [SEQ=0x25, LEN=7, door=0, barrier=0, mode=VENTA, tx=0xFFFFFFFF]} (2.8,-1.6);
    \node[seqnote] at (6.3,-2.2) {ESP32 actualiza lectura de sensores físicos sin requerir ACK adicional};
\end{tikzpicture}
\end{center}

\subsection{Secuencia 6: Diagnóstico de enlace mediante ping/eco (HEARTBEAT)}

\begin{center}
\begin{tikzpicture}
    \node[seqactor] at (2.8,0) (esp) {ESP32};
    \node[seqactor] at (9.8,0) (mega) {Arduino Mega};
    \draw[thick, dashed, gray!60] (2.8,-0.3) -- (2.8,-2.4);
    \draw[thick, dashed, gray!60] (9.8,-0.3) -- (9.8,-2.4);

    \draw[seqmsg] (2.8,-0.7) -- node[above] {HEARTBEAT [SEQ=0x30, LEN=0]} (9.8,-0.7);
    \draw[seqmsgback] (9.8,-1.4) -- node[above] {HEARTBEAT eco [SEQ=0x30, LEN=0]} (2.8,-1.4);
    \node[seqnote] at (6.3,-2.0) {Conectividad básica del enlace serie confirmada};
\end{tikzpicture}
\end{center}


% ============================================================================
% SECCION 9 - EJEMPLOS COMPACTOS DE TRAMAS
% ============================================================================

\clearpage
\section{Ejemplos compactos de tramas}

Los siguientes ejemplos ilustran la composición secuencial byte a byte de tramas
representativas sin campos de verificación checksum o CRC.

\subsection{Ejemplo A: Notificación de tecla \texttt{KEY('5')} emitida por Arduino Mega}

\begin{itemize}[leftmargin=1.6em]
    \item \textbf{Parámetros:} $\texttt{CMD} = \texttt{0x10}$, $\texttt{SEQ} = \texttt{0x14}$, $\texttt{LEN} = 2$, Carácter ASCII \texttt{'5'} (\texttt{0x35}), $\texttt{key\_seq} = \texttt{0x07}$.
    \item \textbf{Trama total transmitida (8 bytes):} \texttt{A5 5A 02 10 14 02 35 07}
\end{itemize}

\noindent\begin{xltabular}{\textwidth}{|p{1.4cm}|p{2.2cm}|p{2.5cm}|X|}
    \hline
    \tablehead{Offset} & \tablehead{Valor (Hex)} & \tablehead{Campo} & \tablehead{Interpretación} \\
    \hline \endhead

    0 & \texttt{0xA5} & \texttt{SOF1} & Primer byte de sincronización. \\
    \hline
    1 & \texttt{0x5A} & \texttt{SOF2} & Segundo byte de sincronización. \\
    \hline
    2 & \texttt{0x02} & \texttt{PROTO} & Versión mayor del protocolo (v2). \\
    \hline
    3 & \texttt{0x10} & \texttt{CMD} & Identificador de comando \texttt{KEY}. \\
    \hline
    4 & \texttt{0x14} & \texttt{SEQ} & Número de secuencia \texttt{0x14} (decimal 20). \\
    \hline
    5 & \texttt{0x02} & \texttt{LEN} & Dos bytes de carga útil. \\
    \hline
    6 & \texttt{0x35} & \texttt{key\_ascii} & Código ASCII del carácter \texttt{'5'}. \\
    \hline
    7 & \texttt{0x07} & \texttt{key\_seq} & Contador secuencial de pulsación local \texttt{0x07}. \\
    \hline
\end{xltabular}

\subsection{Ejemplo B: Orden de dispensado \texttt{VEND(ch=3, tx=42)} emitida por ESP32}

\begin{itemize}[leftmargin=1.6em]
    \item \textbf{Parámetros:} $\texttt{CMD} = \texttt{0x12}$, $\texttt{SEQ} = \texttt{0x20}$, $\texttt{LEN} = 5$, $\texttt{channel} = 3$, $\texttt{transaction\_id} = 42$ (\texttt{0x0000002A}).
    \item \textbf{Trama total transmitida (11 bytes):} \texttt{A5 5A 02 12 20 05 03 2A 00 00 00}
\end{itemize}

\noindent\begin{xltabular}{\textwidth}{|p{1.4cm}|p{2.2cm}|p{2.5cm}|X|}
    \hline
    \tablehead{Offset} & \tablehead{Valor (Hex)} & \tablehead{Campo} & \tablehead{Interpretación} \\
    \hline \endhead

    0 & \texttt{0xA5} & \texttt{SOF1} & Primer byte de sincronización. \\
    \hline
    1 & \texttt{0x5A} & \texttt{SOF2} & Segundo byte de sincronización. \\
    \hline
    2 & \texttt{0x02} & \texttt{PROTO} & Versión mayor del protocolo (v2). \\
    \hline
    3 & \texttt{0x12} & \texttt{CMD} & Identificador de comando \texttt{VEND}. \\
    \hline
    4 & \texttt{0x20} & \texttt{SEQ} & Número de secuencia \texttt{0x20} (decimal 32). \\
    \hline
    5 & \texttt{0x05} & \texttt{LEN} & Cinco bytes de carga útil. \\
    \hline
    6 & \texttt{0x03} & \texttt{channel} & Canal electromecánico solicitado (Canal 3). \\
    \hline
    7 & \texttt{0x2A} & \texttt{tx\_id[0]} & Byte de menor peso de \texttt{transaction\_id} (42 decimal). \\
    \hline
    8 & \texttt{0x00} & \texttt{tx\_id[1]} & Byte 1 de \texttt{transaction\_id}. \\
    \hline
    9 & \texttt{0x00} & \texttt{tx\_id[2]} & Byte 2 de \texttt{transaction\_id}. \\
    \hline
    10 & \texttt{0x00} & \texttt{tx\_id[3]} & Byte de mayor peso de \texttt{transaction\_id}. \\
    \hline
\end{xltabular}


% ============================================================================
% SECCION 10 - CRITERIOS DE VERIFICACION DE INTEGRACION
% ============================================================================

\clearpage
\section{Criterios de verificación de integración}

Las implementaciones de firmware de ESP32 y Arduino Mega deben superar
satisfactoriamente el siguiente banco de pruebas de interoperabilidad antes de
autorizar pruebas en la máquina física:

\noindent\begin{xltabular}{\textwidth}{|p{2.0cm}|p{6.0cm}|X|}
    \hline
    \tablehead{ID Prueba} & \tablehead{Procedimiento de prueba} & \tablehead{Criterio de aceptación normativo} \\
    \hline \endhead

    \texttt{UART-T01} & Inyección de trama sintáctica válida (\texttt{VEND}, \texttt{KEY}, \texttt{DISPLAY}). & El nodo receptor procesa la trama íntegramente y emite la confirmación esperada sin desbordamiento de memoria ni bloqueos. \\
    \hline
    \texttt{UART-T02} & Inyección de trama con valor de longitud anómalo ($\texttt{LEN} > 32$ o no coincidente con \texttt{CMD}). & El receptor descarta la trama por longitud inválida o responde \texttt{ACK\_REJECTED} con motivo \texttt{INVALID\_LENGTH}. \\
    \hline
    \texttt{UART-T03} & Envío de comando \texttt{VEND} indicando un canal fuera de rango ($\texttt{channel} = 0$ o $\texttt{channel} = 5$). & Arduino Mega responde de inmediato con \texttt{ACK\_REJECTED} y motivo \texttt{INVALID\_CHANNEL} sin activar actuadores. \\
    \hline
    \texttt{UART-T04} & Envío sucesivo de dos comandos \texttt{VEND} con idéntico \texttt{transaction\_id}. & Arduino Mega activa el motor una sola vez para la primera orden; ante la segunda devuelve \texttt{ACK\_RECEIVED} sin segundo movimiento mecánico. \\
    \hline
    \texttt{UART-T05} & Envío de comando \texttt{VEND} con idéntico \texttt{transaction\_id} pero canal diferente. & Arduino Mega detecta conflicto de duplicado y rechaza de inmediato con \texttt{ACK\_REJECTED} y motivo \texttt{DUPLICATE\_CONFLICT}. \\
    \hline
    \texttt{UART-T06} & Inyección de trama con identificador de comando no registrado en el diccionario. & El receptor responde de inmediato con \texttt{ACK\_REJECTED} reportando motivo \texttt{INVALID\_CMD}. \\
    \hline
\end{xltabular}

\end{document}
